\section{断面積と体積}\label{節:断面積と体積}

立体を $x$ 軸に垂直な平面で切り,位置 $x$ における断面積を $A(x)$ とする.幅 $\Delta x$ の薄い部分を,底面積 $A(x)$,高さ $\Delta x$ の柱体で近似すれば,その体積は $A(x)\Delta x$ である.断面積が連続に変化する通常の立体では,分割を細かくした極限から
\begin{equation}\label{式:断面積の体積}
V=\int_a^b A(x)\dd x
\end{equation}
を得る.断面が途中で形を変える場合も,変わる位置で有限個の区間に分け,各区間で同じ考えを使える.

この式の $A(x)$ は長さではなく面積である.長さの単位を $\mathrm{cm}$ とすれば,\,$A(x)$ は $\mathrm{cm}^2$,幅 $\dd x$ を掛けた体積は $\mathrm{cm}^3$ になる.図形の式が合っているか迷ったとき,単位も確認の助けになる.

\begin{example}{正方形の断面}
平面上で $0\leq x\leq1$,\,$x^2\leq y\leq x$ となる部分を底面とする.各 $x$ において,底面内の線分を一辺とする正方形が,底面に垂直な断面になるような立体の体積を求めよ.
\end{example}

底面内の線分の長さは $x-x^2$ である.この長さをそのまま積分すると底面積になるが,今求めたいのは体積なので,正方形の断面積
\[
A(x)=(x-x^2)^2=x^2-2x^3+x^4
\]
を積分する.
\[
V=\int_0^1(x^2-2x^3+x^4)\dd x
 =\left[\frac{x^3}{3}-\frac{x^4}{2}+\frac{x^5}{5}\right]_0^1
 =\frac1{30}.
\]
同じ底面でも,断面が正三角形なら面積は $\frac{\sqrt3}{4}(x-x^2)^2$ であり,体積も上の値の $\frac{\sqrt3}{4}$ 倍になる.底面の情報だけでは体積は決まらず,断面の形が必要である.


